\section*{Kindness, Practiced Quietly}
\addcontentsline{toc}{section}{Kindness, Practiced Quietly}

Lurine arrived at the murch on foot, not because she meant to, but because that was how it worked out. The bus stopped short of the outer fence, and the last mile was given over to animals and people who had paid for animals: cows with braided tails and polished horns, carts rented by the hour, boys shouting prices for shade and water and commemorative prints. The air rang with bells and recorded hymns and argument. Reverence had learned how to sell itself.

She stood for a time at the edge of it, pack slung over one shoulder, feeling older than she expected to feel. Not bitter---just tired, in the way one becomes tired of insisting. She had once believed that kindness could win an argument. She no longer believed that. She still believed it mattered.

The murch dominated the far wall, enormous and inevitable. People queued before it in orderly lines, their faces lifted, some weeping, some bored, some intent on the instructions whispered into their ears by guides in gray. Lurine waited her turn without hurry.

When she finally stood before the image, she did not gasp. She did not laugh. She felt, instead, the peculiar tightening that comes when something is both ridiculous and dangerous at once. A face magnified until it became law. A likeness burdened with commandments it had never asked for.

She did not ask herself whom it resembled.
She did not ask herself why.

A young Servant guide stood beside her, reciting softly to a cluster of visitors. He was earnest, barely grown, his robe too large for him.

``Do you accept the Visage as true?'' he asked her, by rote.

Lurine looked at him---not at the image, but at the boy. ``I accept that you're afraid,'' she said.

His mouth opened, closed. ``It is illegal to dispute it.''

``Then don't dispute it,'' she said mildly. ``Just don't let it make you cruel.''

A child tugged at her sleeve then, a pilgrim girl with dust on her cheeks and a paper token tied around her wrist.

``Did you come to see God?'' the child asked.

Lurine smiled, small and honest. ``I came to see what people do when they think they have.''

The air inside the murch felt different than she remembered air feeling anywhere. Less static. More warmth. As if the world were listening back, just a little. She hated that she noticed. She hated more that she liked it.

I like kind people, she remembered thinking once, long ago.
It felt now like both prophecy and accusation.

Near the side aisle, a man knelt with his head bowed, one ear crudely bandaged, blood already seeping through the cloth. A dissenter, probably. Someone who had said one sentence too many. The guards nearby were bored; the punishment had been administered and recorded. There was nothing left to see.

Lurine knelt beside him without asking permission. She took water from her pack, cleaned the wound properly, replaced the bandage with one of her own. She pressed bread into his hand. Not much. Enough.

``Eat slowly,'' she said.

He looked at her as if she might vanish. She did not stay long enough to test the theory.

When she stood and turned away, no one stopped her. No one praised her. The image remained, immense and untroubled, absorbing devotion and coin alike.

Lurine walked out past the cows and the shouting vendors and the bells. She did not feel victorious. She did not feel defeated.

A face could be a prison.
So could a story.

She walked out anyway.

